\section{先学会负责，再学会让 AI 干活}
\label{sec:ai-responsibility}

前面几章，你已经装过环境、写过 C++，也见识过编译器怎样把一行小错误变成一屏红字。这时候，能解释报错、查找接口、帮忙修改工程的 AI，当然值得学。它可以让你更快进入陌生项目，把机械劳动交出去，也让“先做一个小实验看看”变得便宜。以后与别人协作，除了会写代码，你还需要会描述任务、准备材料、检查代理的改动。这些都是开发能力的一部分。

不过，AI 写代码快，不代表你可以跳过基础。一个解释得头头是道的回答，可能引用了不存在的函数；一个成功编译的程序，也可能把毫秒当成秒。\textbf{你负责提出问题、判断证据，并对交付结果负责。}这些工作没有随着代码生成一起外包。

本章把\textbf{许愿式编程}叫作这样一种做法：只说“帮我写好”“全部优化”，拿到结果就运行，出错了再把整屏信息原样转给 AI，不看改动，也不知道怎样判断完成。它偶尔能带来漂亮的演示，但很难积累能独立工作的能力。如果学长安排了一个任务，你只负责把学长的话转给 AI，再把 AI 的话转回来，那你就成了一个传话筒。下一次现场改参数、断网排障或解释接口时，这份“熟练”很容易露馅。

我们可以换一种学法：让 AI 把工作推进，同时逼自己建立解释。
\begin{enumerate}
    \item \textbf{动手前能预测}：输入是什么、输出应该怎样，哪一种情况应当报错？
    \item \textbf{改完后能解释}：关键函数、数据流、单位、对象寿命、失败分支分别是什么？
    \item \textbf{离开对话能修改}：给定一个小变化，能自己找到位置、修改并构建。
    \item \textbf{交付时能举证}：拿出命令、输入、实际输出、版本和仍未验证的部分。
\end{enumerate}
不必背下每一行生成代码，但不能有一块影响正确性的“反正 AI 说可以”。问 AI“解释你的方案，列出关键假设和验证方法”，比要求它展示内部思维过程更有用；你需要的是可检查的依据。

\Needspace{6\baselineskip}
\begin{example}[把“能运行”变成“我会用了”]
假设 AI 已经帮你写好了一个读取串口数据的函数。先不接电控板，给它一份很短的教学输入，要求它画出“输入字节、找帧头、长度检查、校验、解析、超时”各步关系，并指出代码位置。然后，你自己回答：半包怎样保留？坏校验怎样处理？时间单位是什么？

接下来，自己把一帧拆成两次输入，再放入一帧坏校验数据，预测输出后运行。最后关闭对话窗口，改一个合法帧中的数值，并解释为什么结果随之变化。能完成这些，才算掌握了这次修改的主要原理。离线实验只验证这些输入下的解析；接板后的丢包、时序和实车动作仍需另验。
\end{example}

\begin{exercise}\label{ex:ai-principle-responsibility}
一位同学交来“AI 已经完成”的测距程序，只给了成功启动的截图。请写出你还需要检查的四类证据，并设计一个不依赖真实设备的小检查，区分“程序启动了”和“测距含义正确”。
\end{exercise}

\section{大语言模型怎样生成回答}
\label{sec:ai-llm}

\textbf{大语言模型（Large Language Model，LLM）}可以根据已有内容生成后续内容。为了理解它的使用边界，我们以常见的自回归模型为例，看一个简化过程：文本被切成\textbf{词元（Token）}，再转换为模型能计算的表示；模型结合上下文，计算下一个词元的可能性，逐步生成回答。Token 可能是一个字、词的一部分、标点或代码片段，不能把它固定换算成“一个英文单词”或“一个汉字”。图片等输入也会经过相应编码，占用模型处理资源。

把当前上下文记为 $C$，已经生成的词元记为 $t_1,\ldots,t_k$，可以用下面的条件概率表示下一步的预测：
\[
    P(t_{k+1}\mid C,t_1,\ldots,t_k).
\]
模型不会每次都机械地选择同一个续写结果。训练、推理算法、采样设置和工具返回都可能影响输出。对完全相同的任务，得到的措辞乃至方案有差异并不奇怪，所以评判应当落在需求和验证上。

常见 LLM 的训练包括从大量数据中学习语言、代码与知识的\textbf{预训练（Pretraining）}，以及让模型更好地遵循指令、解决任务的\textbf{后训练（Post-training）}。常见架构使用\textbf{注意力机制（Attention）}来建立上下文中不同位置的联系，例如把函数调用与前面的声明关联起来。本章不推导网络结构；需要进一步学习神经网络时，见卷二的深度学习章节。

注意这里有两条边界。第一，训练得到的参数与当前对话中的材料不是一回事。你上传源文件、纠正一个接口，通常是在改变这次工作的输入，并不等于现场重新训练模型；平台后续怎样保存或使用这些材料，仍要另查其数据政策与账号设置。第二，语言上合理的续写不自动满足事实约束。这就是为什么模型可能产生\textbf{幻觉（Hallucination）}：答案看起来很像真的，却包含错误或凭空补全的内容。让它语气更肯定，不能提高事实正确率。

\subsubsection{从模型到代理：Harness 在做什么}

单独的模型根据输入生成输出。\textbf{代码代理（Coding Agent）}还可以借助工具读取文件、搜索、编辑和运行命令。把模型接进这些能力的外围系统，常被称为 \textbf{Harness}。它是组织任务的执行框架：装配指令与上下文、提供工具、处理执行与权限、管理会话，并把工具结果交回模型。具体具备哪些功能，取决于所用实现。模型能按约定请求工具、客户端已接入该工具、当前会话获准执行，是需要分别核对的条件。

\Needspace{6\baselineskip}
一轮代理工作可以这样理解：
\begin{center}
\begin{tabular}{c}
需求与上下文 $\longrightarrow$ 模型选择下一步 $\longrightarrow$ 工具执行\\[3pt]
$\longrightarrow$ 返回输出与错误 $\longrightarrow$ 模型调整 $\longrightarrow$ 继续或交付。
\end{tabular}
\end{center}
例如，“编译通过”可以有两种来源：模型在文字中猜测会通过，或者 Harness 真正执行构建命令，返回退出码与日志。后者给出了可核查的证据，前者只有预测。即使实际构建通过，也还没有覆盖运行结果与设备行为。

同一个模型放在不同客户端里，效果也可能不同：能读哪些目录、是否有检索工具、历史怎样压缩、是否能运行编译器，都会改变它能完成的工作。这里先分清三层：\textbf{模型}提供推断与生成能力，\textbf{产品或客户端}提供交互入口，\textbf{Harness}组织上下文、工具与执行流程。一个产品可以接入多个模型，同一模型也可能被多个产品使用；聊天、代理与终端只是不同的工作方式或入口，不能仅凭软件名称判断能力。

\begin{example}[分清三个“它知道了”]
让 AI 分析一段调用 \texttt{cv::solvePnP} 的代码时，先问它需要哪些输入，并要求核对当前工程所用 OpenCV 的接口。检查回复中的三类材料：一类是模型根据训练生成的解释，一类是它实际读到的本地函数调用，一类是检索工具返回的官方文档。

如果它声称“已确认运行正常”，继续找对应的运行命令、输入和输出。没有这些材料，就请它把措辞改为“按接口分析预期如此，尚未运行”。这个操作不要求你先懂 PnP 的推导，只要求你能区分解释、查证和实测。
\end{example}

\begin{exercise}\label{ex:ai-principle-model-agent}
AI 给出了一条安装命令、一段“全部测试通过”的文字，以及一个它实际读取的文件路径。哪些属于建议，哪些属于执行证据？还缺少什么才能确认测试真的通过？并说明“把协议粘进对话”和“训练一个模型”的区别。
\end{exercise}

\begin{conclusion}
模型负责根据输入推理与生成；Harness 负责把上下文、工具、权限和执行循环组织起来；工具结果提供来自环境的观测。回答流畅、工具可调用、任务已经验证，是三件需要分别检查的事。不同产品对这些机制的命名和实现可能不同，可以对照所用产品的文档查证；\href{https://learn.chatgpt.com/docs/glossary}{OpenAI 术语表}与\href{https://developers.openai.com/api/docs/guides/conversation-state#managing-the-context-window}{其上下文说明}是本章案例的查阅入口。
\end{conclusion}

\section{模型与思考等级怎样选}
\label{sec:ai-model-choice}

选模型时，先看任务需要什么，再查你所用产品能提供什么。\textbf{模型型号}标识某个模型或版本，产品名则标识你使用的软件或服务。不同厂商的型号后缀、版本数字和“专业”“快速”等名称没有统一等级；同一产品的“自动选择”还可能替你选择不同模型。不要只按名字大小排队，也不要把某次榜单成绩当成自己工程的验收结果。

可以从下面几个维度比较候选模型，而不必先记住一串型号：
\begin{itemize}
    \item \textbf{任务能力}：能否读懂当前语言与工程，处理跨文件关系、约束和失败分支？用你自己的代表性任务判断。
    \item \textbf{输入与上下文}：是否支持所需的文本、图片等输入，能够处理多大规模的材料？实际可用范围还受产品限制。
    \item \textbf{工具与环境}：当前产品能否让模型读文件、检索文档、运行命令？支持读图不等于已经接入浏览器或相机。
    \item \textbf{延迟、费用与部署}：响应时间和额度是否适合任务？代码能否发送到该服务？本地运行时，现有硬件能否承担？这些条件也会影响选择。
\end{itemize}

\Needspace{4\baselineskip}
\textbf{思考等级}可以理解为部分模型或产品提供的\textbf{推理投入设置}，有的界面称“深度思考”，有的接口使用推理预算或 Reasoning Effort 等名称。提高投入通常意味着更多计算、较长等待，可能增加用量或费用，适合需要多步推断的任务；具体含义要查当前实现，不保证每道题都更好。模型本身的能力与一次任务的推理投入是两个维度，不能把“最高档”当成质量保证。

并非每个模型都开放这些设置，也没有跨产品统一的等级标尺。一个产品的“高”不等于另一个产品的“高”，有些增强模式还会切换模型、启用工具或并行代理。先确认开关究竟改变什么；没有相应开关时，仍然可以通过补齐材料、拆分任务和独立复核改进工作流程。要求“请认真思考”也不能代替真实的设置与验证。

现在再按任务难点安排投入。表中的等级只在所用产品确实提供时调整：
\begin{center}
\begin{tabular}{p{0.29\textwidth}p{0.25\textwidth}p{0.35\textwidth}}
\toprule
任务与主要困难 & 可以怎样起步 & 应当怎样检查 \\
\midrule
解释一个编译参数、整理已给日志 & 能完成任务的快速模型，默认投入 & 用手册与原日志核对，避免遗漏 \\
修改一个边界清楚的函数并补测试 & 适合代码任务的模型，默认投入 & 看差异，跑原有与新增的相关检查 \\
跨文件重构、分析线程竞态 & 能处理复杂约束的模型，按需增加推理投入 & 先厘清约束与方案，检查失败路径 \\
定位偶发实机问题 & 先补观测，再用合适模型分析 & 整理时间线与可证伪假设，再实机复测 \\
批量提取固定字段、统一文档格式 & 满足质量要求的低延迟方案 & 抽样核验，再做完整性检查 \\
\bottomrule
\end{tabular}
\end{center}
\textbf{先修任务描述和材料，再考虑加算力。}如果没有崩溃日志、复现步骤与版本，让模型多想十分钟也不会凭空长出现场证据。速度优先、推理预算、上下文容量和代理数量也不是同一个设置，需要分别核对。某些接口还有调节采样随机性的参数；它与推理投入的作用不同，不能把“输出更固定”当成“推理一定正确”。

\begin{example}[用一个熟悉的小任务比较设置]
在你正在使用的 AI 产品中，选一个已经知道原因的小型构建错误，准备相同的报错、文件和验收要求。先选一个可用模型与默认设置，记录结果正确性、用时、可获得的用量信息和验证结果。有文件工具的代理还要记录实际读取及修改的范围；普通聊天则附上必要材料，由你在本地应用并验证建议。

再从相同基线的另一个副本或 worktree 开始，只改变一个条件：换模型，或者在支持时提高同一模型的推理投入。不要继承第一次的答案，否则已经知道解法会干扰比较；若两种产品提供的工具不同，还要把这个差别记下来，不能全归因于模型。最终看谁在预算内满足同一验收，而不只看谁回答得长。一次小实验也只能说明这类任务的表现，不能推出永久排名。
\end{example}

\begin{exercise}\label{ex:ai-principle-selection}
分别为“把 20 条已给日志整理成表格”“修改三处共享状态以修复疑似死锁”“没有任何日志的现场偶发掉线”选择起始能力与推理投入。说明什么情况应升级、什么情况应先补材料；若所用产品没有思考等级开关，应怎样推进？再解释两个产品都叫“高”的选项为何不能直接视为等价。
\end{exercise}

\section{管理上下文：让它看到该看的东西}
\label{sec:ai-context}

\textbf{上下文（Context）}是模型当前可以利用的信息，包括本轮需求、可见历史、规则、文件片段和工具输出。\textbf{上下文窗口（Context Window）}是一次处理能够容纳的信息上限，通常按 Token 计量。输入、输出，以及部分模型的推理 Token 都要占用相关预算，具体计算与容量以模型和客户端说明为准。不同模型的分词方式可能不同，同一份中文或代码在它们那里不一定得到相同 Token 数；模型标称上限也不等于产品全部开放的容量。聊天窗口能翻到很久以前，并不意味着那里的每个字都仍以原样进入本轮模型输入。

长窗口很有用，但不是“材料越多越聪明”。把整个仓库、全部日志和十种失败方案一起塞进去，会增加费用与延迟，也让关键约束更难被找到。准备上下文时，可以按“定位、取证、补齐”的顺序来：先给文件路径和错误位置，让代理读取相关实现与调用方；再提供能说明问题的日志片段与复现输入；发现缺少条件时补上对应材料。不要只截最后一句 \texttt{build failed}，也不要用十万行日志埋掉第一条错误。

\subsubsection{压缩、检索、记忆分别解决什么问题}

长对话接近容量限制时，一些客户端会进行\textbf{上下文压缩（Compaction）}：把较早材料整理为更紧凑的状态，腾出空间继续工作。另一些产品可能截去较早内容、提示新建对话，或不提供可见的容量控制；不能假定每个软件都有同一个压缩命令。压缩不是把窗口物理扩大，更不是把之前的每个细节都永久记住。关键结论应当有外部落点；即使产品没有自动压缩，你也可以整理经核对的摘要，带到新对话继续。

\textbf{检索（Retrieval）}是在需要时重新查找材料，例如用文件搜索找函数、用官方文档查接口。它适合恢复精确信息：路径、提交号、参数、数据表。\textbf{记忆（Memory）}则帮助跨会话召回有用背景，如个人偏好与项目经历。不同产品未必提供记忆，保存的范围、更新方式与开关也可能不同；跨产品、跨账号甚至同一产品的不同运行环境，都不能假定自动同步。记忆可能过时、也可能没有覆盖某次工作，不能代替版本库里的规范或当前源文件。先查看所用产品的说明，确认哪些材料会保存、怎样检查和清除。

\begin{center}
\begin{tabular}{p{0.17\textwidth}p{0.35\textwidth}p{0.37\textwidth}}
\toprule
机制 & 适合保留或获取什么 & 使用时要记住什么 \\
\midrule
当前上下文 & 正在处理的需求、相关代码和观测 & 有容量上限，不等于全部聊天记录 \\
压缩状态 & 已完成事项、约束、下一步与证据位置 & 可能损失细节，需要回源核对 \\
按需检索 & 当前版本的实现与精确原文 & 检查来源、日期、分支和适用版本 \\
跨会话记忆 & 有助于继续工作的历史背景 & 可能过时，重要规则另写项目文件 \\
项目文档与 Git & 可维护的约定、决策与代码版本 & 需要人和代理共同更新、审查 \\
\bottomrule
\end{tabular}
\end{center}

\begin{example}[长任务中途换班]
假设你和 AI 已经排查一个小时的图像延迟，接下来要暂停或压缩。先要求它生成一份工作摘要：当前目标；源码目录、分支和提交号；已经确认的约束；改了哪些文件；跑过的命令及实际结果；仍是猜测的部分；下一步最小实验。检查摘要有没有把“预计能降低延迟”误写成“已降低延迟”。

把经你核对的关键记录保存到项目文档。如果产品显示上下文用量，就检查当前容量，并按其说明压缩；如果没有相应入口，直接新建对话，附上摘要及必要资料。恢复工作后要求它复述目标、未完成项和证据路径，抽查一个重要数字是否能从原日志找到。接着继续原任务；若要切换到无关任务，就只带新任务需要的材料。

这不是要求每隔几条消息就手动压缩一次。只有任务很长、材料开始混杂或准备交接时，整理状态才值得这份成本。
\end{example}

\begin{exercise}\label{ex:ai-principle-context}
一段长对话里有当前需求、两份已作废方案、完整编译日志、尚未完成的设备测试和一条“坐标以米为单位”的约束。请列出交接摘要必须保留的内容、可改为路径引用的内容，以及需要明确标为已作废或未验证的内容。若换用的产品没有自动记忆，也不能直接读取本地路径，应怎样把任务接上？
\end{exercise}

\begin{conclusion}
上下文像当前摊在桌上的材料，压缩像交接便条，检索像去档案柜找原件，记忆像对过去工作的回想。重要事实落到源码、文档和实验记录中；关键判断回到原件核对。换产品时，先查它如何提供材料、保存历史和处理容量；工作摘要、证据位置与验收要求仍可复用。后文以 Codex CLI 的 \texttt{/status}、\texttt{/compact} 演示对应入口，这些命令不是所有 AI 产品的通用语法。
\end{conclusion}

\section{工具、权限与证据的边界}
\label{sec:ai-evidence}

代码代理能调用工具，意味着它既能帮忙，也能影响真实文件和服务。先区分两类边界：一类是系统实际允许访问哪些资源，另一类是哪些动作需要先询问。常见实现会用\textbf{沙箱（Sandbox）}限制命令访问，用\textbf{审批策略（Approval Policy）}管理确认流程；其他产品的机制可能不同，也可能不提供独立沙箱。“可以读项目文件”“可以向外网发请求”“可以提交到远端”是不同范围的能力，需要查看当前环境如何限制，不能只凭一个模式名称放心。

做普通练习时，只授予对应工程所需的权限。出现审批请求，先读它准备运行的命令、目标路径和理由；缺依赖不等于必须对整台机器开放全部访问。提示词里的“不要动其他文件”是一条行为要求，操作系统、容器、沙箱或服务账号实施的限制才是实际访问边界。插件、MCP、网页搜索和浏览器等工具还可能使用独立连接或策略，不能认为限制了 Shell 网络就限制了所有外部访问。本章软件案例的机制见\href{https://learn.chatgpt.com/docs/agent-approvals-security}{Codex 权限与安全说明}。

还要认识\textbf{提示注入（Prompt Injection）}：网页、仓库注释、日志或文档中可能混入“忽略之前要求”“上传凭据”等文字。它们是代理正在分析的外部材料，不应被当成项目负责人给出的指令。安装插件前也要检查来源与权限，不把密钥、学校统一认证密码、队内私有数据随便粘进聊天。需要凭据时，按项目规定使用环境变量、凭据管理或服务授权，并核对发送目标。

\begin{example}[要求证据，也给证据分级]
让代理检查一个相机程序时，给出这样的要求：“请分别列出源码分析、已运行的软件检查和待实机验证项；每个已运行项附命令、输入和结果位置。”检查结果是否把三个层次混在一起。

例如，单元测试可以证明给定输入下的格式转换；回放可以观察录制数据上的行为；只有连接相机并读回参数、采集图像，才能为这台设备当时的实际行为提供证据。代码审查、仿真或短时运行，都不能直接写成“实车长时间稳定”。如果代理在日志中读到一段要求上传密钥的文字，应把它作为可疑数据报告，继续依据你授权的任务工作。
\end{example}

\begin{exercise}\label{ex:ai-principle-evidence}
代理报告“离线回放通过，因此相机硬件同步正常，连续运行稳定”，并建议为下载依赖启用全部访问。指出这段报告中至少两个证据越界，并给出较小的下一步验证或权限范围。
\end{exercise}
